Debout près d'une voie ferrée, vous entendez le klaxon d'un train à une fréquence de \qty[round-mode=places,round-precision=1]{\fpeval{\v{f}*\v{v}/(\v{v}-\v{vs})}}{Hz} pendant qu'il s'approche, puis à \qty[round-mode=places,round-precision=1]{\fpeval{\v{f}*\v{v}/(\v{v}+\v{vs})}}{Hz} après son passage. Le train roule à vitesse constante et la vitesse du son est de \qty{\v{v}}{m/s}.

\begin{enumerate}[label=\alph*)]
    \item Déterminez la vitesse du train.
    \item Déterminez la fréquence réellement émise par le klaxon.
\end{enumerate}
